\documentclass[../../../../main.tex]{subfiles}
\begin{document}

\begin{flushright}
\footnotesize\textit{【自動文字起こし・要確認】}
\end{flushright}

{\bf [解]}
$\triangle BCD$の外心を$H$とおく．$H$を通り，$\triangle BCD$に垂直な直線$\ell$上の点を$X$とする．この時，外心の定義から
\begin{align*}
\overline{XB}=\overline{XC}=\overline{XD}\quad\cdots①
\end{align*}
である．$AB$の中点を$E$とし，$E$を通り$AB$に垂直な平面を$\pi$とする．$AB\perp\pi$，$\ell\perp$平面$BCD$から，$\pi$と$\ell$は必ず交点を持ち，これを$P$とすると，
\begin{align*}
\overline{AP}=\overline{BP}\quad\cdots②
\end{align*}

①，②から，$P$を中心とし，半径$\overline{AP}$の球面は$A,B,C,D$を全て通るので，題意は示された．$\blacksquare$

\begin{figure}[htb]
\centering
\begin{tikzpicture}[scale=1.4, >=stealth]
  \coordinate (A) at (0,2.6);
  \coordinate (B) at (-2,0);
  \coordinate (C) at (0.6,-0.5);
  \coordinate (D) at (2,0.5);
  \coordinate (H) at (barycentric cs:B=1,C=1,D=1);
  \coordinate (P) at ($(A)!0.45!(H)$);
  \draw (A) -- (B) -- (C) -- (D) -- cycle;
  \draw (A) -- (C);
  \draw[dashed] (B) -- (D);
  \draw[dashed] (A) -- (H);
  \draw[dashed] (A) -- (P) -- (B);
  \draw[dashed] (P) -- (C);
  \draw[dashed] (P) -- (D);
  \node[above] at (A) {$A$};
  \node[left] at (B) {$B$};
  \node[below] at (C) {$C$};
  \node[right] at (D) {$D$};
  \node[below right] at (H) {$H$};
  \node[above right] at (P) {$P$};
  \fill (H) circle (0.6pt);
  \fill (P) circle (0.6pt);
\end{tikzpicture}
\caption{四面体$ABCD$と$\triangle BCD$の外心$H$，外接球の中心$P$}
\end{figure}

\newpage

\end{document}
